% Transcription — fonds Alexandre Grothendieck, shelfmark 151, batch 4 (pages 61-74).
% Established from the facsimile archives/batches/151/batch-04.pdf (a mirror of
% https://grothendieck.umontpellier.fr/151.pdf), per the skill
% transcribe-grothendieck.
%
% This is the last batch of the folder and it is fourteen pages, not twenty:
% the PDF holds 75 pages, one of them Montpellier's cover sheet, so the
% archivists' numbering stops at 74. All fourteen carry the mathematics; none
% is skipped.
%
% The author numbers his sheets 19 to 25 at the head of pages 62, 64, 66, 68,
% 69, 71, 73 — continuing the series batch 3 carried to 18. The parity of the
% numbered side flips between sheets 22 and 23 (page 68 carries 22, page 69
% carries 23), so from page 69 on the number stands at the head of the odd
% pages. This matters for one cross-reference: on page 71 he sends the reader
% to « le th. énoncé p. 22 », which is the theorem of recollement stated on
% archivists' page 68.
%
% The batch falls in two parts. Pages 61-68 finish section 3, « Stratifications
% globales : introduction des tubes » — the fundamental diagram (17), the
% census of which inclusions actually hold between the elementary pieces
% X_i^*, V^*_{ij}, V_{ij}, the diagram Ĩ built on the ordered set I with its
% three types of arrow, worked examples for Delta_r and for two small posets,
% the dimension function d : I -> Z, and on page 68 a first statement of the
% « Th. de recollement des tubes ». Pages 69-74 open section 4, « Topos
% canoniques associés à une stratification globale », whose equation numbering
% restarts: A) general inverse images, B) the case of an X_{I'} attached to a
% crible I' of I, C) the V^{I'''}_{I',I''}, which the folder breaks off in the
% middle of.
%
% Two heavily reworked passages are transcribed as skeletons and marked as
% such with \note: the struck paragraph on page 68, crossed through with two
% long diagonal strokes, and the opening of page 70, which a long diagonal
% marginal note runs across. Pages 61, 68, 70 and 71 carry marginal notes
% written diagonally up the left edge; where they do not read we say so rather
% than offering a reading.
%
% On uncertainty: \ill marks what we do not read; \uncertain marks a reading
% that carries mathematical weight and that we do not hold to be sure.
% Function words whose sense is forced by the sentence are given plainly.
%
% Two numbering decisions worth recording. The three inclusion types displayed
% on page 63 carry manuscript labels we cannot read, so they are given
% unnumbered and a \note says so. The block of axioms opening section 4 on page
% 69 is read as (1) — the glyph would also admit (11), but the equations that
% follow it run (2) to (11) across pages 69-71, and page 71's « qui vient de
% (1 d) » points back to its part d).
%
% The dating is the inventory's own; the notes themselves carry no date.
%
% Pass: Opus 5 (claude-opus-5), 2026-08-16 — first pass, unchecked against the pages by a human.
\documentclass[11pt,a4paper]{article}
% Le préambule commun à toutes les transcriptions du fonds Grothendieck.
%
% Il tient en une idée : **l'incertitude se compose, elle ne se résout pas.**
% Une transcription qui lisse un mot douteux a détruit la seule chose pour
% laquelle on la faisait — un lecteur doit pouvoir savoir, à chaque endroit,
% ce qui était sur le feuillet et ce qui a été ajouté. Les six macros qui
% suivent sont donc l'appareil critique, et non de la décoration.
%
% Les mêmes commandes sont reconnues par `scripts/render.mjs`, qui produit la
% vue de lecture du site. Une macro ajoutée ici doit l'être aussi là-bas,
% faute de quoi le rendu échouera bruyamment — ce qui est voulu : un rendu
% silencieusement faux est pire qu'un rendu absent.

\ProvidesPackage{grothendieck}[2026/08/08 Transcriptions du fonds Grothendieck]

\usepackage{amsmath,amssymb,amsthm}
\usepackage{mathrsfs}
\usepackage{stmaryrd}
% Les diagrammes commutatifs sont partout dans ce fonds : c'est la langue
% même de la géométrie algébrique de Grothendieck, et une transcription qui
% les rendrait en prose ne transcrirait rien.
\usepackage{tikz-cd}
% Français par défaut — c'est la langue des feuillets — mais l'anglais est
% chargé aussi : la traduction et le résumé sont des documents anglais, et la
% typographie française (espaces avant les deux-points) y serait une faute.
% Ces documents basculent par \selectlanguage{english} en tête de corps.
\usepackage[english,french]{babel}
\usepackage{fontspec}
\usepackage{geometry}
\geometry{a4paper,margin=25mm,marginparwidth=28mm}
\usepackage{marginnote}
\usepackage{xcolor}
% ulem, not soul. `soul`'s \st and \ul are built on a character-by-character
% scan that XeLaTeX's font machinery breaks: the document fails with an
% undefined control sequence rather than degrading. ulem does the same two jobs
% and is Unicode-engine safe. `normalem` stops it hijacking \emph, which the
% transcriptions use for Grothendieck's own emphasis.
\usepackage[normalem]{ulem}
\usepackage{hyperref}
\hypersetup{colorlinks=true,linkcolor=black,urlcolor=black,pdfborder={0 0 0}}

% --- Métadonnées du lot -----------------------------------------------------
% Reprises telles quelles par le script de rendu, qui les affiche en tête de
% la vue de lecture. Elles fixent ce que le fichier prétend couvrir : sans
% elles, une transcription flotte sans point d'attache dans un fonds de seize
% mille feuillets.

\newcommand{\folder}[1]{\gdef\grothFolder{#1}}
\newcommand{\batch}[1]{\gdef\grothBatch{#1}}
\newcommand{\leaves}[2]{\gdef\grothFirst{#1}\gdef\grothLast{#2}}
\newcommand{\foldertitle}[1]{\gdef\grothFoldertitle{#1}}
\gdef\grothFolder{?}\gdef\grothBatch{?}\gdef\grothFirst{?}\gdef\grothLast{?}\gdef\grothFoldertitle{}

% --- Le résumé --------------------------------------------------------------
%
% Il ouvre la lecture modernisée, sous le titre « Résumé », et il est composé
% en petit. Ce n'est pas un ornement : c'est le seul passage du document qui
% ne soit pas des mathématiques, et un lecteur doit pouvoir le reconnaître
% avant de l'avoir lu — puis le sauter s'il connaît déjà le sujet, ou s'y
% arrêter s'il ne le connaît pas. Le filet qui le suit marque la reprise du
% corps.

\newenvironment{resume}
  {\small\setlength{\parskip}{0.45em}}
  {\par\medskip\hrule\medskip}

% --- Mots-clés ---------------------------------------------------------------
%
% En anglais, en clôture du résumé de la lecture modernisée : le vocabulaire
% moderne sous lequel chercher ce que les feuillets construisent. Le manifeste
% du site (`npm run manifest`) les extrait pour étiqueter la cote entière —
% cette ligne est la source unique des tags, il n'y a pas de fichier à côté.
\newcommand{\keywords}[1]{\par\smallskip\noindent
  {\footnotesize\textsc{Keywords} --- \textit{#1}}\par}

% --- L'appareil critique ----------------------------------------------------

% Le numéro de feuillet, en marge. C'est l'ancre qui permet de régler un
% désaccord sur une formule en regardant *un* feuillet plutôt qu'une chemise
% de six cents. Le site s'en sert aussi pour tourner les pages du fac-similé
% quand on fait défiler la transcription.
\newcommand{\leaf}[1]{%
  \marginnote{\footnotesize\color{gray!70}\textsf{#1}}%
  \label{leaf:#1}\ignorespaces}

% Illisible : on ne devine pas. Un mot inventé qui se lit comme les autres est
% le pire résultat possible.
\newcommand{\ill}{\textcolor{red!60!black}{[\,\ldots\,]}}

% Lecture douteuse : le mot est donné, et signalé comme incertain.
\newcommand{\uncertain}[1]{\uline{#1}}

% Ajout de l'éditeur — une lettre manquante, un mot sous-entendu.
\newcommand{\add}[1]{\textcolor{blue!50!black}{[#1]}}

% Biffé par Grothendieck lui-même. Ce qu'il a rayé fait partie du document :
% une rature dit souvent où la pensée a changé de direction.
\newcommand{\struck}[1]{\sout{#1}}

% Note du transcripteur, sur l'état matériel du feuillet.
\newcommand{\note}[1]{\par\smallskip{\small\color{gray!60!black}\textit{[#1]}}\par\smallskip}

% Note marginale de Grothendieck, distincte du corps du texte.
\newcommand{\marginal}[1]{\par\smallskip\noindent\hspace{1em}%
  {\small\color{gray!60!black}#1}\par\smallskip}

% --- Titre ------------------------------------------------------------------

\AtBeginDocument{%
  \begin{center}
    {\large\bfseries Fonds Alexandre Grothendieck --- cote n\textsuperscript{o}~\grothFolder}\\[2pt]
    {\itshape\grothFoldertitle}\\[4pt]
    {\small feuillets \grothFirst--\grothLast\ (lot \grothBatch)}\\[2pt]
    {\footnotesize\color{gray!60!black}Transcription établie d'après le fac-similé numérisé
     par l'Université de Montpellier.\\
     Les lectures incertaines sont soulignées, les illisibles notées [\,\ldots\,],
     les ajouts entre crochets.}
  \end{center}
  \vspace{1em}}

\endinput


\folder{151}
\batch{4}
\pages{61}{74}
\dating{[à partir de 1981-à partir de 1990]}
\watermark{Édition de démonstration}
\foldertitle{Topos stratifié (1981) : notes manuscrites (s.d.).}

\begin{document}

\section*{Introduction des tubes — suite}

\page{61}
Posons déjà

\[
  (14) \qquad
  \begin{cases}
    \mathcal{V}_{\Delta_1} = \coprod\limits_{\substack{(i,j) \in I^2 \\ i \leq j}}
      \mathcal{V}_{i,j} \\[10pt]
    \mathcal{V}^{*}_{\Delta_1}
      = \coprod\limits_{\substack{(i,j) \in I^2 \\ i \leq j}} \mathcal{V}^{*}_{i,j}
      \qquad \subset \mathcal{V}_{\Delta_1}
  \end{cases}
\]

\note{la seconde ligne s'ouvre sur un « $\mathcal{V}$ » biffé ; le $\subset$
porte « ouvert » écrit dessous}

on a donc

\[
  (15)
\]

\begin{tikzcd}
  \mathcal{V}^{*}_{\Delta_1} \arrow[r, hook] \arrow[d]
    & \mathcal{V}_{\Delta_1} \arrow[d, dashed]
    & X^{*}_{\Delta_0} \arrow[l] \arrow[dl, dashed] \\
  X^{*}_{\Delta_0} & X &
\end{tikzcd}

\struck{\ill} qui reflètent les diagrammes

\[
  (17)
\]

\begin{tikzcd}
  \mathcal{V}^{*}_{i,j} \arrow[r] \arrow[d]
    & \mathcal{V}_{i,j} \arrow[d, dashed]
    & X^{*}_{i} \arrow[l] \arrow[dl, dashed] \\
  X^{*}_{j} & X &
\end{tikzcd}

\marginal{NB Le \ill}\note{la note marginale, écrite en diagonale au bas à
gauche, accompagne un petit diagramme ordonné — un sommet d'où partent deux
traits, l'un vers un sommet qui en domine un troisième ; nous ne la relions pas
avec certitude au texte}

\struck{Notons} Finalement, on n'est pas sûr que les notations $(14)$ soient
bien utiles — du fait \struck{\ill} que le carré dans $(15)$ n'a rien de
cocartésien, car le diagramme en trait plein, qui est le diagramme « sous-

\page{62}\note{p. 19 de l'auteur}
jacent » des diagrammes $(17)$, définit une \emph{limite} qui se reproduit dans
\ill\ comme espace topologique, a peu de chances d'être homéo. à $X$ (p.ex. si
$X$ connexe \ldots) !

Les $\mathcal{V}^{\,j}_{i,k}$ sont des germes de sous-espaces de $X$, dont les
seuls \struck{relations d\ill} \struck{morphismes} canoniques entre eux sont
les morphismes d'inclusion — à quelle condition a-t-on

\[
  (18) \qquad \mathcal{V}^{\,j}_{i,k} \subset \mathcal{V}^{\,j'}_{i',k'}
\]
\[
  (19) \qquad i \leq j \leq k , \qquad i' \leq j' \leq k'
\]

donnés ? Question \ill\ il \ill\ pour les espaces élémentaires

\[
  X_i^{*} , \quad \mathcal{V}^{*}_{i,j} , \quad \mathcal{V}_{i,j}
  \qquad (j \text{ successeur de } i) \ ?
\]

\struck{Cela} Je vais \struck{\ill} supposer pour simplifier
\[
  \begin{cases}
    X_i^{*} \neq \emptyset \\
    X_i \subset X_j \Rightarrow i \leq j
  \end{cases}
\]
— car sinon on aurait $X_i^{*} = \emptyset$ (justifie \ldots) Il n'y a donc pas
de relations d'inclusion entre les $X_i^{*}$ pour des indices distincts.
\emph{Si} \struck{pour} $\mathcal{V}_{i,j} \subset \mathcal{V}_{i',j'}$,
prenons deux « pts » \struck{des $X''$} \ill\ genres d'après, il faut avoir
$X_i^{*} \subset X_{i'}^{*}$ d'où $i = i'$ ; \ill\ chose de \ill\ que les
\uncertain{$j$ le long des $X_i^{*}$} déterminent $X_j$ ($j$ successeur de $i$)
\ill\ bien voisinages \ill\ $\mathcal{V}$) \ill\ déjà $j$ ce qui est \ill\ devra
avoir aussi $j = j'$. De même, on voit que si
$\mathcal{V}^{*}_{i,j} \subset \mathcal{V}^{*}_{i',j'}$, \ill\ un des pts
d'adhérence, \struck{\ill} ce qui met $\overline{X_i^{*}}$ et $X_{i'}$ \ill\
\struck{$X_{i'}$ et $X_{i'}$}, sont en relation d'inclusion,

\note{la page est ici surchargée de reprises ; nous donnons le fil de
l'argument et laissons en \ill\ ce qui ne se lit pas}

\page{63}
dans le cas (raisonnable) où $\overline{X_i^{*}} = X_i$ $\forall i$ ; on \ill\
donc $i \leq i'$. \struck{Dans le cas} \ill\ les cas $i = i'$, on voit comme
tout à l'heure que \ill\ (« raisonnable ») $j = j'$. Mais le cas $i < i'$
est-il possible ? Comme
$\emptyset \neq \mathcal{V}^{*}_{i,j} \subset X_j^{*}$,
$\mathcal{V}^{*}_{i',j'} \subset X_{j'}^{*}$, on doit avoir
$X_j^{*} \cap X_{j'}^{*} \neq \emptyset$ donc $j = j'$ si l'on avait \ill\
(ci !)

\begin{tikzcd}[column sep=small, row sep=small, nodes={font=\scriptsize}]
  & i' \arrow[r, no head] & j' & & X_{i'} \arrow[r, no head] & X_{j'} \\
  i \arrow[ur, no head] \arrow[dr, no head] & & & X_i \arrow[ur, no head] \arrow[dr, no head] & & \\
  & j & & & X_j &
\end{tikzcd}

donc $i \leq i' < j = j'$ et par suite $i' = i$ puisque $j$ successeur de $i$.

On voit de même que les seuls morphismes d'inclusion d'un $X_i^{*}$ dans un des
autres types est

\[
  X_i^{*} \longrightarrow \mathcal{V}_{i,j}
  \qquad (\text{dans les cas « raisonnables »})
\]

\struck{il n'y a pas de \ill\ morphismes d'un $\mathcal{V}_{i,j}$ dans un des
autres types}, \struck{et $\mathcal{V}_{i,j}$} enfin que le seul morphisme d'un
$\mathcal{V}^{*}_{i,j}$ dans un des autres types est

\[
  \mathcal{V}^{*}_{i,j} \longrightarrow X_j^{*}
\]
\[
  \mathcal{V}^{*}_{i,j} \longrightarrow \mathcal{V}_{i,j}
\]

\note{les trois inclusions portent au manuscrit des numéros que nous ne lisons
pas}

en définitive, il n'y a plus lieu d'envisager d'autres morphismes entre ces
types — pièces \struck{élémentaires}.

\page{64}\note{p. 20 de l'auteur}
Ce \struck{sont} aussi les \struck{diagrammes} flèches en traits pleins des
diagrammes fondamentaux $(17)$.

Avec l'ens ordonné $I$, fabriquons donc un \struck{ens. ordonné}
\emph{diagramme} dont les sommets sont \struck{considérés \ill} \ill\ des
symboles $[i]$, $[i,j]$, $[i,j]^{*}$, où $i,j \in I$, $j$ successeur de $i$.

Les flèches du diagramme sont de trois types, qui sont toutes décrites dans

\[
  (19)
\]

\begin{tikzcd}
  {[i,j]}^{*} \arrow[r] \arrow[d] & {[i,j]} & i^{*} \arrow[l] \\
  {[j]}^{*} & &
\end{tikzcd}

Désignons par $\widetilde{I}$ ce diagramme, dont l'ens. des sommets
\struck{s'identifie à} \struck{\ill} \marginal{en cardinal} card
$I + 2\mathcal{V}$, \struck{\ill} où $\mathcal{V}$ \marginal{card} est l'ens.
des drapeaux $(i<j)$ « élémentaires », \marginal{i.e. tels que $j$ soit
successeur de $i$} et de \ill\ des flèches est $3\mathcal{V}$. Si $I$ contient
exactement deux éléments $i,j$ avec $i<j$, \marginal{si $I \simeq \Delta_1$} (ce
diagramme est celui de $(19)$). Si $I = \Delta_r = (0,1,\ldots,r)$, on trouve
(en explicitant maintenant les symboles par les $X_i^{*}$,
$\mathcal{V}^{*}_{ij}$, $\mathcal{V}_{ij}$ qu'ils désignent) le diagramme

\page{65}
\struck{\ill}\note{la page s'ouvre sur un début de tracé rayé, puis sur une
première mise en place du même escalier, portée en haut à droite et annotée
« NB on \ill » ; nous ne donnons que la version au net}

\[
  (20) \qquad \widetilde{\Delta}_n
\]

\begin{tikzcd}[column sep=small, row sep=small, nodes={font=\scriptsize}]
  X_0^{*} \arrow[r] & \mathcal{V}_{0,1} & \mathcal{V}^{*}_{0,1} \arrow[l] \arrow[d] & & & & & & \\
  & & X_1^{*} \arrow[r] & \mathcal{V}_{1,2} & \mathcal{V}^{*}_{1,2} \arrow[l] \arrow[d] & & & & \\
  & & & & X_2^{*} \arrow[r] & \mathcal{V}_{2,3} & \mathcal{V}^{*}_{2,3} \arrow[l] \arrow[d] & & \\
  & & & & & & X_3^{*} & & \\
  & & & & & & \cdots \arrow[d] & & \\
  & & & & & & X^{*}_{n-1} \arrow[r] & \mathcal{V}_{n-1,n} & \mathcal{V}^{*}_{n-1,n} \arrow[l] \arrow[d] \\
  & & & & & & & & X_n^{*}
\end{tikzcd}

Si $I$ est le diagramme $a \leftarrow \bullet \rightarrow b$, on trouve

\[
  (21)
\]

\begin{tikzcd}[column sep=small, nodes={font=\scriptsize}]
  \mathcal{V}^{*}_{0,b} \arrow[r] \arrow[d] & \mathcal{V}_{0,b}
    & X_0^{*} \arrow[l] \arrow[r] & \mathcal{V}_{0,a}
    & \mathcal{V}^{*}_{0,a} \arrow[l] \arrow[d] \\
  X_b^{*} & & & & X_a^{*}
\end{tikzcd}

Si $I$ est le diagramme $a \rightarrow c \leftarrow b$, on trouve

\[
  (22)
\]

\begin{tikzcd}
  X_a^{*} \arrow[d] & & X_b^{*} \arrow[d] \\
  \mathcal{V}_{a,c} & & \mathcal{V}_{b,c} \\
  \mathcal{V}^{*}_{a,c} \arrow[u] \arrow[r] & X_c^{*} & \mathcal{V}^{*}_{b,c} \arrow[u] \arrow[l]
\end{tikzcd}

etc.

\struck{La première chose qu'il faudrait \ill\ préciser est le} Supposons qu'on
dispose sur $I$ d'une fonction à valeurs entières

\page{66}\note{p. 21 de l'auteur}
\[
  (23) \qquad I \xrightarrow{\ d\ } \mathbb{Z}
  \qquad \Bigl( \text{i.e.} \quad I = \coprod_m I_m \Bigr)
\]

jouant le rôle d'une fonction dimension (on \ill\ une présentation en quelque
sorte \ill\ avec une fonction type codimension), satisfaisant la condition

\[
  (24) \qquad \text{si } i,j \in I ,\ i \leq j ,\
  \text{ alors } d(j) = d(i)+1 \iff j \text{ succes. de } i
\]

\note{« $I$ » est suivi d'un $\emptyset$ biffé}

Posons alors, pour $m, n \in \mathbb{Z}$

\[
  (25) \qquad
  \begin{cases}
    X_m^{*} = \coprod\limits_{i \in I_m} X_i^{*} \\[10pt]
    \mathcal{V}_{m,m+1} = \coprod\limits_{\substack{(i,j) \in I \times I \\
      i \leq j ,\ d(i)=m ,\ d(j)=m+1}} \mathcal{V}_{i,j} \\[16pt]
    \mathcal{V}^{*}_{m,m+1} = \coprod\limits_{i,j \text{ comme dessus}}
      \mathcal{V}^{*}_{i,j}
  \end{cases}
\]

de sorte qu'on a les diagrammes

\[
  (26) \qquad (\mathrm{Diag})_m
\]

\begin{tikzcd}
  X_m^{*} \arrow[r] & \mathcal{V}_{m,m+1} & \mathcal{V}^{*}_{m,m+1} \arrow[l] \arrow[d] \\
  & & X^{*}_{m+1}
\end{tikzcd}

s'enchaînant dans un diagramme comme

Enfin en principe, mais que nous \struck{\ill} c'est-à-dire \ill\ forme
\struck{$d(I)d$} \ill\ simplifier $I$ fini, et en n'écrivant pas les
\struck{parties du diagramme} \struck{\ill} correspondant à des \struck{\ill}
sommets occupés par des espaces vides :

\note{les deux dernières lignes de la page sont reprises trois fois et
biffées ; nous n'en donnons que ce qui se lit}

\page{67}
Dans le cas d'une fonction \emph{codimension} ($\underline{d(j) = j-1}$ si $j$
succ. de $i$), on \struck{pourra} considérera des $\mathcal{V}_{d,d-1}$,
$\mathcal{V}^{*}_{d,d-1}$ au lieu des $\mathcal{V}_{d,d+1}$,
$\mathcal{V}^{*}_{d,d+1}$, et le diagramme prend la forme

\note{le manuscrit écrit bien « $j-1$ » ; le parallèle avec $(24)$ demanderait
« $d(i)-1$ »}

\struck{(27)} (cas d'une fonction codimension)

\[
  (27)
\]

\begin{tikzcd}[column sep=small, row sep=small, nodes={font=\scriptsize}]
  X_n^{*} \arrow[r] & \mathcal{V}_{n,n-1} & \mathcal{V}^{*}_{n,n-1} \arrow[l] \arrow[d] & & & & \\
  & & X^{*}_{n-1} \arrow[r] & \mathcal{V}_{n-1,n-2} & \mathcal{V}^{*}_{n-1,n-2} \arrow[l] \arrow[d] & & \\
  & & & & X^{*}_{n-2} & & \\
  & & & & \cdots \arrow[d] & & \\
  & & & & X_1^{*} \arrow[r] & \mathcal{V}_{1,0} & \mathcal{V}^{*}_{1,0} \arrow[l] \arrow[d] \\
  & & & & & & X_0^{*}
\end{tikzcd}

\page{68}\note{p. 22 de l'auteur — c'est la page que la référence « le th.
énoncé p. 22 » de la page 71 désigne}
\textbf{Th. de recollement des tubes} (première version) Considérons le
diagramme de type $\widetilde{I}$, formé par les $X_i^{*}$,
$\mathcal{V}^{*}_{i,j}$, $\mathcal{V}_{i,j}$ ($i,j \in I$, $j$ successeur de
$i$) et les trois types de flèches \marginal{précédents} y figurant (en trait
plein) \struck{explicités} et indiqué\add{s} dans $(17)$. Alors $X$ s'identifie
\struck{au topos} à la \ill\ de ce type.

\marginal{pour que ce soit précis, \ill\ formes \ill\ $X$ \ill\ déterminé par
$\mathrm{Top}(X)$ \ldots}\note{note marginale écrite en diagonale le long du
bord gauche, reprise plusieurs fois ; nous n'en donnons que ce qui se lit}

\struck{Faute de disposer de définitions \ill\ pour les voisinages tubulaires,
je vais faire une démonstration, lorsqu'on \ill\ en remplaçant
$\mathcal{V}_{i,j}$ par des \marginal{vrais} \ill\ des
$\mathcal{V}_{i,j} \setminus R_{i,j} \cap \mathcal{V}_{i,j}$, où $X_j$, de
façon que $\mathcal{V}_{i,j}$ est un voisinage de $X_i$ dans $X_j$, et qu'on
pose $\mathcal{V}^{*}_{i,j} = \mathcal{V}_{i,j} \setminus X_i^{*}$.
\struck{Alors}}

\note{tout ce paragraphe est barré de deux longs traits diagonaux ; nous en
donnons le squelette}

\struck{démonstration} On va en donner une \ill\ heuristique ci-dessous. Mais
avant de le faire, et vu la généralisation, on va \ill\ un énoncé qui permet
d'exprimer ci-\ill\ les « topos canoniques » associés à ce \ill\ en termes
\ill\ de type $\widetilde{I}$ formé des topos élémentaires.

\section*{4. Topos canoniques associés à une stratification globale}

\page{69}\note{p. 23 de l'auteur}
On va montrer comment, à une situation stratifiée donnée, on peut en associer
d'autres.

\medskip

A) \emph{Images inverses générales.} Rappelons les axiomes utilisés jusqu'à
présent :

\[
  (1) \qquad
  \begin{cases}
    \text{a)} & X_i \text{ fermés},\ (X_i)_{i \in I} \text{ loc. fini} \\
    \text{b)} & i \leq j \Rightarrow X_i \subset X_j \\
    \text{c)} & X = \bigcup X_i \\
    \text{d)} & \forall i,j \in I ,\ \ldots \quad
      X_i \cap X_j = \bigcup\limits_{k \leq i,j} X_k
  \end{cases}
\]

Notons que pour tout $X'$ au-dessus de $X$, la famille des

\[
  (2) \qquad X_i' \overset{\text{déf}}{=} X_i \times_X X'
\]

\note{le signe $=$ est suivi d'un mot biffé que nous ne lisons pas}

satisfait les \ill\ mêmes conditions.

D'ailleurs le système $\tfrac{1}{2}$ simpl. des $X_{\Delta_r}$ — comme image
inverse \ill\ des $X'_{\Delta_r}$, défini par les $X_i'$, \struck{Ain}si on a

\[
  (3) \qquad
  \begin{cases}
    X'_{\Delta_r} \simeq X_{\Delta_r} \times_X X' \\
    X'^{\,*}_{\Delta_r} \simeq X^{*}_{\Delta_r} \times_X X' \\
    \dot{X}'_{\Delta_r} = \dot{X}_{\Delta_r} \times_X X'
  \end{cases}
\]

\page{70}
\marginal{NB} Nous appliquerons ces \ill\ surtout au cas où $X'$ est un ouvert
de $X$. C'est pour pouvoir prendre de telles images inverses \ill, qu'il
n'\ill\ pas été commode \ill, de supposer les $X_i$ ou les $X_i^{*}$ non-vides,
\struck{ou encore} que $I \mapsto X_i$ est un plongement d'ens ordonnés
$I \hookrightarrow \mathcal{P}(X)$.

Lorsque $X' \to X$ est une immersion locale \struck{mais pas} (\ill\ n'est une
immersion ouverte !) alors \ill\ \struck{la définition} les images inverses
\struck{\ill} \ill\ fournies de $X$ comment : la formation des voisinages
tubulaires d\add{e} une telle partie dans une entière. Notons

\note{le haut de cette page est traversé par une longue note marginale
diagonale, plusieurs fois reprise, qui recouvre le texte ; nous en donnons le
fil et laissons en \ill\ ce qu'elle rend illisible}

d'ailleurs que pour $i \leq j$, on aura (sous l'hyp. d'ailleurs que $X' \to X$
soit une immersion locale)

\[
  (4) \qquad
  \begin{cases}
    R'_{ij} = R_{ij} \times_X X' \\
    S'_{ij} = S_{ij} \times_X X'
  \end{cases}
\]

d'où, en le cas d'une immersion locale \marginal{propre}, des iso

\[
  (5) \qquad
  \begin{cases}
    \mathcal{V}'_{i,j} \simeq \mathcal{V}_{i,j} \times_X X' \\
    \mathcal{V}'^{\,*}_{i,j} \simeq \mathcal{V}^{*}_{i,j} \times_X X'
  \end{cases}
  \qquad i \leq j
\]

\page{71}\note{p. 24 de l'auteur}
et plus généralement

\[
  (6) \qquad \mathcal{V}'^{\,j}_{i,k} \simeq \mathcal{V}^{\,j}_{i,k} \times_X X' ,
\]

\ill.

Ceci montre en particulier que la démonstration du th. de recollement, i.e. du
th. énoncé p. 22, est une \ill\ question \emph{locale} sur $X$ — ce qui permet
p.ex. \ill\ vides — on \ill\ si \emph{$I$ est fini}.

\marginal{\ill}\note{la marge gauche porte ici une longue note diagonale et un
appel de note $\circledast$ ; ni l'une ni l'autre ne se lisent}

\struck{Considérons p.ex. le cas où $X' = $}

\medskip

B) \emph{Cas d'un $X_{I'}$.} \struck{Soit} Soit $I'$ une partie de $I$ telle que

\[
  (7) \qquad i \leq j \in I' \implies i \in I'
\]

et soit

\[
  (8) \qquad X_{I'} = \bigcup_{i \in I'} X_i \qquad (\text{partie fermée de } X)
\]

\[
  (9) \qquad \text{On a bien sûr } I' \subset I'' \implies X_{I'} \subset X_{I''} ,
  \text{ plus gén.!}
\]
\[
  (10) \qquad X_{I' \cup I''} = X_{I'} \cup X_{I''}
  \qquad \text{plus gén.}
  \begin{cases}
    X_{\cup I'_\alpha} = \bigcup\limits_\alpha X_{I'_\alpha} \\
    \text{avec } X_\emptyset = \emptyset
  \end{cases}
\]

et aussi

\[
  (11) \qquad X_{I' \cap I''} = X_{I'} \cap X_{I''}
  \qquad (\text{donc } X_I = X)
\]

qui vient de $(1\ \mathrm{d})$. Dans ces formules,

\page{72}
$I'$, $I''$, les $I'_\alpha$ sont des parties de $I$ satisfaisant $(7)$ (des
\emph{cribles} de $I$)

Si dans A) on prend $X' = X$, il est plus commode de travailler avec la
stratification de $X'$ définie par les $X_i$ avec $i \in I'$ — il est clair que
les conditions $(1)$ \struck{relatives} : $X' = X_{I'}$ sont satisfaites. Les
« parties cribles » de $X'$ pour cette stratification, on \ill\ celles induites
au sens général des espaces au-dessus de $X$, sont les mêmes — ou surtout les
traces \marginal{au} \struck{induites} sur $X' = X_{I'}$ des parties-cribles de
l'espace stratifié $X$.

Ici, les espaces élémentaires pour la stratification de type $I'$ de
$X' = X_{I'}$, sont les espaces

\[
  (12) \qquad X_i^{*} ,\ \mathcal{V}^{*}_{ij} ,\ \mathcal{V}_{ij}
  \quad \text{pour } i,j \in I' .
\]

\struck{Combinant les deux opérations, on a p.ex. une stratification pour la
\ill\ $\mathcal{U}'$ d'un $X' = X_{I'}$, on a les espaces \struck{élémentaires}
associés, qui sont donc les}

\[
  (13) \qquad \mathcal{U}' \cap X_i^{*} ,\
  \mathcal{U}' \cap \mathcal{V}^{*}_{ij} ,\
  \mathcal{U} \cap \mathcal{V}_{ij} \qquad (i,j \in I')
\]

\note{ce dernier paragraphe et la formule $(13)$ sont barrés d'une grande
croix ; la formule est donnée ici sans rature, la rature portant sur le bloc
entier}

\page{73}\note{p. 25 de l'auteur}
Revenons pour un instant à $X$, et considérons l'ens \struck{des} $I_0$ des
$i \in I$ tels que $X_i = \emptyset$. C'est un crible, et on a
$X_i^{*} = \emptyset$, $\mathcal{V}^{*}_{ij} = \mathcal{V}_{ij} = \emptyset$ si
$i \in I_0$. Donc négligeant les espaces vides, on voit que le diagramme de
type $\widetilde{I}$ défini par l'espace stratifié $X$ ne change pas quand on
\ill\ son \ill\ $I$ par $I \setminus I_0$, ou plus gén. par
$I \setminus I'_0$, où $I'_0 \subset I_0$ \ill\ est un crible, \ill\ ce qui
donne lieu : un diagramme $\widetilde{I \setminus I'_0}$ qui est
\emph{contenu} dans $\widetilde{I}$ (cela est vrai pour \emph{tout} crible de
$I$).

Si p.ex. on a deux cribles

\[
  (14) \qquad I'' \subset I' \subset I
\]

d'où

\[
  (15) \qquad X_{I''} \subset X_{I'} \quad \text{et} \quad
  \mathcal{U}_{I',I''} \overset{\text{déf}}{=} X_{I'} \setminus X_{I''} ,
\]

on peut \struck{\ill} (et on a) sur $\mathcal{U}_{I',I''}$ regarder plutôt la
\struck{filtration} stratification de type $I' \setminus I''$, définie par les
\struck{$X_{i'} \setminus X_{i'} \cap X_{I''}$ ($i' \in I'$) et}

\[
  (16) \qquad
  X_{i',\, \mathcal{U}_{I',I''}}
  = X_{i'} \setminus X_{i'} \cap X_{I''} \subset \mathcal{U}_{I',I''}
\]

dont les topos élémentaires sont donc les $X_{i'}^{*}$
($i' \in I' \setminus I''$) et des $\mathcal{V}^{*}_{i,j',X'}$,
$\mathcal{V}_{i,j',X'}$

\page{74}
associés aux couples $(i',j')$ avec $i' \in I' \setminus I''$, $j'$ successeur
de $i'$. On a

\[
  (17)
\]

\begin{tikzcd}
  \mathcal{V}_{i',j',X'} \arrow[r, hook] & \mathcal{V}_{i',j'} \\
  \mathcal{V}^{*}_{i',j',X'} \arrow[r, hook] \arrow[u, hook]
    & \mathcal{V}^{*}_{i',j'} \arrow[u, hook]
\end{tikzcd}

mais il n'est pas clair en général que ce soient \struck{des}
\marginal{seulement} (des) équivalences d'homotopie \ldots

\struck{Donc} il faut mieux ne pas trop s'\ill\ quand il s'agit de comparer les
constructions sur un $X_{I'}$, et sur un
$\mathcal{U}_{I',I''} = X_{I'} \setminus X_{I''}$.

Je vais en disposer par la suite des \struck{\ill} plus particulières ou qui se
\ill\ pour ce s'introduisent ainsi sur un ouvert $\mathcal{U}_{I',I''} \ldots$

\medskip

C) \emph{Les $\mathcal{V}^{\,I'''}_{I',I''}$.} On se donne trois cribles

\[
  (18) \qquad
  \begin{array}{c} I' \\ I''' \end{array} \subset I'' \subset I ,
\]

d'où

\[
  (19) \qquad
  \begin{array}{c} X_{I'} \\ X_{I'''} \end{array}
  \subset X_{I''} \hookrightarrow X_I = X
\]

\end{document}
